\documentclass[a4,11pt]{aleph-notas}

% -- Paquetes adicionales
\usepackage{enumitem}
\usepackage{aleph-comandos}
\usepackage{systeme}

% -- Datos  
\institucion{Facultad de Ciencias Exactas, Naturales y Ambientales}
\carrera{Ciencia de datos}
\asignatura{Álgebra lineal}
\tema{Material extra no. 5: Errores frecuentes - Matrices}
\autor{Andrés Merino}
\fecha{Periodo 2026-1}

\logouno[0.16\textwidth]{Logos/logoPUCE_04_ac}
\definecolor{colortext}{HTML}{1957d1}
\definecolor{colordef}{HTML}{1957d1}
\fuente{montserrat}

% -- Comandos adicionales
\setlist[enumerate]{label=\roman*.}

\begin{document}

\encabezado

\section{Resolución de ejercicios sin procedimiento}

\subsection*{Incorrecto:}
No escribir nada en la resolución:
\[
    \begin{pmatrix}
    x \\ y \\ z
    \end{pmatrix}
    =
    \begin{pmatrix}
    -1 & 1 & 1\\
    6 & -5 & -3\\
    -2 & 2 & 1
    \end{pmatrix}
    \begin{pmatrix}
    2 \\ 4 \\ 1
    \end{pmatrix}
    =
    \begin{pmatrix}
    3 \\ -11 \\ 5
    \end{pmatrix}
\]

\subsection*{Correcto:}
Dado que la matriz de coeficientes es invertible, la solución está dada por $A^{-1}b$. Entonces:
\[
    \begin{pmatrix}
    x \\ y \\ z
    \end{pmatrix}
    =
    \begin{pmatrix}
    -1 & 1 & 1\\
    6 & -5 & -3\\
    -2 & 2 & 1
    \end{pmatrix}
    \begin{pmatrix}
    2 \\ 4 \\ 1
    \end{pmatrix}
    =
    \begin{pmatrix}
    3 \\ -11 \\ 5
    \end{pmatrix}.
\]
Por lo tanto,
\[
    x = 3,\quad y = -11,\quad z = 5.
\]

\section{Operaciones por filas no válidas}

\subsection*{Incorrecto:}

Estas operaciones no corresponden a transformaciones elementales válidas:
\begin{itemize}
    \item \(3F_1 + F_2 \to F_3\)
    \item \(F_1 - 3F_2 \to F_2\)
    \item \(F_2 - F_2 \to F_2\)
    \item \(3F_2 \to F_1\)
\end{itemize}


\subsection*{Correcto:}

Las únicas operaciones elementales por filas son:
\begin{enumerate}
    \item Intercambio de filas: \(F_i \leftrightarrow F_j\)
    \item Multiplicación por escalar no nulo: \(cF_i \to F_i\), \(c \neq 0\)
    \item Suma de múltiplos de filas: \(cF_j + F_i\to F_i\)
\end{enumerate}

\section{Uso incorrecto de los símbolos}

\subsection*{Incorrecto:}
Las siguientes matrices no son iguales:
\[
\begin{pmatrix}
1 & 2\\
3 & 4
\end{pmatrix}
=
\begin{pmatrix}
1 & 2\\
0 & -2
\end{pmatrix}
\qquad -3F_1 + F_2 \to F_2.
\]

\subsection*{Correcto:}

Cuando se aplican operaciones por filas, se debe usar $\sim$:
\[
\begin{pmatrix}
1 & 2\\
3 & 4
\end{pmatrix}
\sim
\begin{pmatrix}
1 & 2\\
0 & -2
\end{pmatrix}
\qquad -3F_1 + F_2 \to F_2.
\]

El símbolo \(=\) se usa para igualdad estricta (misma matriz). El símbolo \(\sim\) se usa cuando se tiene equivalencia por filas

\section{Uso incorrecto del símbolo igual}

\subsection*{Incorrecto:}

Colocar igual donde no corresponde, esto es incorrecto porque una matriz no es igual a un escalar.
\[
A =
\begin{pmatrix}
1 & 2\\
3 & 4
\end{pmatrix} = 1\cdot 4 - 2\cdot 3 = -2
\]

\subsection*{Correcto:}
Se debe usar la notación correcta:
\[
\det(A) =
\begin{vmatrix}
1 & 2\\
3 & 4
\end{vmatrix}
= -2
\]



\end{document}